\section{Một corpus đủ nhỏ để chạy, đủ lệch để đo}
\label{sec:ch05-corpus}

\index{corpus đồ chơi!ba ràng buộc}%
Bốn mục sau đều là số đo, nên trước hết phải nói rõ đo trên cái gì. Bài huấn
luyện trên mười hai triệu cặp \tn{câu nguồn}{source sentence} và
\tn{câu đích}{target sentence} của WMT'14 Anh-Pháp, mười ngày trên một máy tám
GPU.
Sách này hứa với bạn rằng mọi thí nghiệm chạy xong trên một laptop không có card
đồ hoạ, nên tôi không đi theo đường đó được, và nói rõ ngay từ đầu rằng cái tôi
dựng là đồ chơi.

Ba ràng buộc phải thoả cùng lúc, và không tập dữ liệu công khai nào thoả cả ba.
Nó phải chạy xong trên CPU trong ngân sách của sách. Nó phải tái lập được từ một
seed, không tải gì, để \code{verify.py} chạy được trên máy người khác. Và hai
thứ tiếng phải khác thứ tự từ đủ nhiều, vì chương~\ref{ch:dong-chinh} cần một
\tn{ma trận đồng chỉnh}{alignment matrix} có chỗ bắt chéo chứ không phải một
đường chéo.

Nên corpus được sinh ra từ một văn phạm nhỏ nằm trong repo. Anh sang Việt, cùng
chiều với bài, vì tiếng Anh là câu nguồn ở cả hai bên.

\subsection*{Thứ duy nhất nó mô hình hoá thật}

\index{cụm danh từ!đảo vị trí tính từ}%
Đó là \tn{cụm danh từ}{noun phrase} đảo, và \tn{đảo câu nguồn}{source reversal}
ở mục~\ref{sec:ch05-dao-cau-nguon} là chuyện khác hẳn, đừng lẫn hai cái. Tiếng
Anh đặt tính từ trước danh từ, tiếng Việt đặt sau. Tiếng Anh đánh dấu xác định
bằng mạo từ, tiếng Việt dùng \tn{loại từ}{classifier}, và loại từ không có từ
tiếng Anh nào tương ứng.

\begin{minted}{text}
the black cat  ->  con  mèo  đen
a  new  lamp   ->  một cái  đèn  mới
\end{minted}

Một mô hình dịch corpus này bằng cách chép vị trí sang vị trí sẽ sai tính từ ở
mọi cụm có tính từ. Và đó là chỗ đáng giá: một ma trận đồng chỉnh đường chéo
không giải được corpus này, nên một mô hình đồng chỉnh được nó thì đã học được
cái gì đó thật.

Văn phạm sinh câu một hoặc hai mệnh đề, mỗi mệnh đề là chủ ngữ, động từ, tân
ngữ, cộng một cụm chỉ nơi chốn tuỳ chọn. Mỗi động từ tiếng Anh đúng một token,
còn phía tiếng Việt thì một hoặc hai (\enquote{muốn} một token,
\enquote{nhìn thấy} hai), nên độ dài câu đích không phải hàm của độ dài câu
nguồn.

\begin{minted}{text}
  en  the white bird wants the chair at the window
  vi  con chim trắng muốn cái ghế cạnh cửa sổ
  en  the cat wants a dog
  vi  con mèo muốn một con chó
\end{minted}

\begin{measured}{corpus sinh từ văn phạm, tag \code{ch05}}
\begin{minted}{text}
split  pairs  src min  src max  src mean  tgt min  tgt max  tgt mean
train  6000   5        21       12.232    5        25       13.716
test   300    5        21       12.397    6        24       13.877

target longer than source in 0.7800 of training pairs
source sentences shared between train and test: 0
\end{minted}

Hai \tn{từ điển}{vocabulary} đều đóng: 34 loại token bên nguồn, 37 bên đích, kể
cả ba ký hiệu đặc biệt. Cả hai dựng từ chính văn phạm chứ không từ một mẫu, nên
số hiệu token không đổi khi thí nghiệm rút nhiều câu hay ít câu.

Đo bằng \code{experiments/ch05\_corpus.py} tại tag \code{ch05}.
\end{measured}

Câu đích dài hơn câu nguồn ở 0.7800 số cặp, và loại từ là lý do chính. Con số
\enquote{0 câu chung giữa train và test} là đo chứ không phải giả định: tập test
được rút dư gấp đôi rồi lọc bỏ mọi câu đã có trong train.

\subsection*{Thứ nó không mô hình hoá, và chương sẽ không giả vờ}

\index{corpus đồ chơi!giới hạn}%
Danh sách này quan trọng ngang danh sách trên, vì mọi con số ở bốn mục sau phải
đọc trong khung của nó.

Văn phạm hữu hạn và phi ngữ cảnh, nên một mô hình đủ lớn học thuộc được nó. Ở
đây không có hình thái, không có hợp giống số, không có nhập nhằng, không có từ
hiếm. Từ điển đóng và mọi token xuất hiện hàng nghìn lần, nên bài toán
\tn{ngoài từ điển}{out-of-vocabulary} không thể xảy ra ở đây. Chính bài toán ấy
ăn mất một phần điểm Bilingual Evaluation Understudy, viết tắt là BLEU, của bài
báo, và bài nhắc tới nó ngay trong tóm tắt. Và câu
thì vô nghĩa về nội dung: văn phạm ràng buộc cú pháp, không ràng buộc nghĩa, nên
\enquote{con mèo mang một cái đèn mới} là một câu hợp lệ của nó.

Vậy nên không con số nào trong chương này so được với BLEU của bài. Chúng đo
những thứ khác nhau trên những tập khác nhau ở những quy mô khác nhau. Cái
chúng làm được là kiểm các \emph{lập luận} của bài, và lập luận thì không đổi
theo quy mô: nếu đảo câu nguồn có tác dụng vì nó rút ngắn khoảng cách giữa từ
tương ứng, thì tác dụng ấy phải thấy được ở mọi quy mô mà khoảng cách còn có
nghĩa.

\subsection*{Cách chấm điểm}

\index{khớp đúng cả câu!vì sao không dùng BLEU}%
Mọi bảng dưới đây chấm \tn{khớp đúng cả câu}{exact match}: cả câu đích đúng
từng token thì được 1, sai một token thì được 0. Không dùng BLEU, có chủ ý. Trên
một corpus nhỏ và đều thế này, một điểm cho phần đúng đo một thứ không đáng tin;
còn cái hỏng mà chương này quan tâm là mô hình đặt tính từ sai bên, tức một câu
sai chứ không phải một câu hơi kém hơn.

Hai cột lặp lại ở mọi bảng. \code{short} là câu nguồn từ 11 token trở xuống,
\code{long} là dài hơn. Ngưỡng 11 là biên giữa một mệnh đề và hai mệnh đề của
văn phạm, chọn trước khi nhìn kết quả chứ không phải sau; trong 300 câu test thì
148 câu ngắn và 152 câu dài. Tách hai cột ấy ra là chuyện đáng làm, và mục sau
là chỗ cho thấy vì sao.

\subsection*{Thiết lập huấn luyện, và ba chỗ nó lệch khỏi bài}

\index{corpus đồ chơi!thiết lập huấn luyện dùng chung}%
Bốn bảng của chương này đặt cạnh nhau được thì chúng phải cùng một thiết lập,
nên thiết lập ấy nằm một chỗ trong \code{seq2seq.py} chứ không chép vào từng
script: 6000 cặp huấn luyện, 300 cặp test, batch 128, 14 epoch, \(d_h = 128\),
embedding 32 chiều.

Ba chỗ lệch khỏi bài, và cả ba là chuyện ngân sách chứ không phải bất đồng.

Optimizer là Adam ở learning rate 0.005. Bài dùng SGD thường ở learning rate
0.7, giảm nửa mỗi nửa epoch sau epoch thứ năm, tổng cộng 7.5 epoch. Bài có mười
ngày trên tám GPU để chờ lịch ấy hội tụ; tôi có vài chục giây.

Cắt gradient thì \emph{không} lệch: ngưỡng norm 5.0, đúng con số bài đặt ở mục
3.4, và chương~\ref{ch:phan-tich} đã dẫn ra vì sao phép cắt ấy hoạt động.

Và 14 epoch là ngân sách chứ không phải hội tụ. Cùng mô hình chạy 20 epoch cho
0.9533 và chạy 40 epoch cho 0.9600, so với 0.8533 ở 14. Tôi in số ở 14 vì mười
lăm lần huấn luyện phải nằm gọn trong ngân sách cả lần chạy, và vì mọi so sánh
trong chương chỉ cần hai vế dừng ở cùng một chỗ. Khi đọc một con số tuyệt đối
trong chương này, nhớ rằng nó còn cách chỗ mô hình đi được khoảng mười điểm.
